% =============================================================================
%  Notas de clase — Sesión 20: Taller de Criptografía
%  Matemáticas de la Computación — Universidad Panamericana
%  Compilar: pdflatex sesion20_taller_criptografia.tex
% =============================================================================
\documentclass[12pt, oneside]{article}
\usepackage{matcomp}
\usepackage{array}

\begin{document}

\portada{Sesión 20}{Actividad: Taller de Criptografía}{2026}

\tableofcontents
\newpage

% =============================================================================
\section{Objetivo}
% =============================================================================

En sesiones anteriores implementaste funciones de codificación y
decodificación para tres sistemas criptográficos: el \textbf{cifrado
por desplazamiento} (\textit{shift cipher}), el \textbf{cifrado afín}
(\textit{affine cipher}) y \textbf{RSA}.  En esta actividad pondrás
esas funciones a trabajar.

Tu equipo elegirá un verso o estrofa de una canción en español y lo
cifrará con los tres sistemas.  Los resultados se entregan a través
del Google Forms indicado al final de la sección~\ref{sec:entrega}.

\begin{recuerda}
  Antes de cifrar, convierte tu frase completa a \textbf{mayúsculas}.
  Los tres sistemas operan exclusivamente sobre los 48~caracteres de
  la tabla de la sección~\ref{sec:tabla}.
\end{recuerda}

% =============================================================================
\section{Tabla de caracteres}
\label{sec:tabla}
% =============================================================================

Todos los cifrados de esta actividad usan la siguiente tabla.  A cada
carácter le corresponde un código decimal de \textbf{dos dígitos}
(de \texttt{00} a \texttt{47}).  El tamaño del alfabeto es $N = 48$.

\begin{center}
\renewcommand{\arraystretch}{1.3}
\begin{tabular}{>{\centering}p{1.0cm} >{\centering}p{0.8cm} |
                >{\centering}p{1.0cm} >{\centering}p{0.8cm} |
                >{\centering}p{1.0cm} >{\centering}p{0.8cm} |
                >{\centering}p{1.0cm} >{\centering}p{0.8cm} |
                >{\centering}p{1.0cm} >{\centering\arraybackslash}p{0.8cm}}
\toprule
\textbf{Car.} & \textbf{Cód.} &
\textbf{Car.} & \textbf{Cód.} &
\textbf{Car.} & \textbf{Cód.} &
\textbf{Car.} & \textbf{Cód.} &
\textbf{Car.} & \textbf{Cód.} \\
\midrule
A & 00 & B & 01 & C & 02 & D & 03 & E & 04 \\
F & 05 & G & 06 & H & 07 & I & 08 & J & 09 \\
K & 10 & L & 11 & M & 12 & N & 13 & O & 14 \\
P & 15 & Q & 16 & R & 17 & S & 18 & T & 19 \\
U & 20 & V & 21 & W & 22 & X & 23 & Y & 24 \\
Z & 25 & Ñ & 26 & 0 & 27 & 1 & 28 & 2 & 29 \\
3 & 30 & 4 & 31 & 5 & 32 & 6 & 33 & 7 & 34 \\
8 & 35 & 9 & 36 & . & 37 & , & 38 & : & 39 \\
\texttt{[esp]} & 40 & ¡ & 41 & ! & 42 & Á & 43 & É & 44 \\
Í & 45 & Ó & 46 & Ú & 47 & & & & \\
\bottomrule
\end{tabular}
\end{center}

\smallskip
\noindent{\small \texttt{[esp]} representa el carácter \textbf{espacio} (código~40).}

% =============================================================================
\section{Instrucciones generales}
% =============================================================================

\subsection{Equipo}

La actividad se realiza en \textbf{equipos de hasta tres personas}.
Todos los integrantes deben escribir sus nombres al responder el Forms.

\subsection{La frase}

Elige un \textbf{verso o estrofa de una canción en español} que cumpla
con todo lo siguiente:

\begin{itemize}
  \item Escrita originalmente en español (no traducción).
  \item Claramente identificable: incluye el título de la canción y
        el artista o grupo.
  \item Sin groserías ni doble sentido ofensivo.
  \item Una vez convertida a mayúsculas, debe contener únicamente
        caracteres de la tabla de la sección~\ref{sec:tabla}.
\end{itemize}

\begin{consejo}
  Elige un verso que tenga entre 20 y 60 caracteres para que el proceso
  sea manejable a mano y el bloqueo RSA sea interesante.
\end{consejo}

\subsection{¿Qué entregar en el Forms?}
\label{sec:entrega}

Deberás repetir el proceso para cada uno de los tres cifrados.  Por
cada cifrado, el Forms pedirá:

\begin{enumerate}
  \item \textbf{Nombre de la canción y artista.}
  \item \textbf{Frase original en mayúsculas} (el texto plano).
  \item \textbf{Clave de codificación} (varía según el cifrado).
  \item \textbf{Frase cifrada} (el texto cifrado).
  \item \textbf{Clave de decodificación} — solo para RSA.
\end{enumerate}

Accede al formulario en:\quad
\url{https://forms.gle/9hBfr1R9P7mp9R1T9}

% =============================================================================
\section{Cifrado por Desplazamiento}
% =============================================================================

El \textbf{cifrado por desplazamiento} suma un número fijo $k$ al código
de cada carácter, módulo $N = 48$.  Para codificar el carácter con
código $m_i$:

\[
  c_i = (m_i + k) \bmod 48
\]

Cada $c_i$ es también el código de un carácter de la tabla; la
frase cifrada es la cadena de esos caracteres.

\begin{description}
  \item[Clave de codificación a reportar.] El entero $k$ que elegiste,
        con $1 \le k \le 47$.
\end{description}

\begin{consejo}
  Para verificar que tu función decodifica correctamente, aplica
  $m_i = (c_i - k + 48) \bmod 48$ y comprueba que recuperas el texto
  original.
\end{consejo}

% =============================================================================
\section{Cifrado Afín}
% =============================================================================

El \textbf{cifrado afín} combina una multiplicación y un desplazamiento:

\[
  c_i = (a \cdot m_i + b) \bmod 48
\]

Para que la función sea invertible, $a$ debe ser coprimo con $N = 48$.

\begin{recuerda}
  Como $48 = 2^4 \times 3$, un valor $a$ es coprimo con 48 si y solo si
  $a$ \textbf{no es divisible} ni por 2 ni por 3.  Por ejemplo: 5, 7,
  11, 13, 17, 19, 23, 25\ldots\ son válidos; 4, 6, 9 no lo son.
  El valor $a = 1$ equivale a un cifrado por desplazamiento.
\end{recuerda}

Al igual que en el cifrado por desplazamiento, cada $c_i$ es el código
de un carácter de la tabla; entrega la frase cifrada como cadena de
caracteres.

\begin{description}
  \item[Clave de codificación a reportar.] El par $(a,\, b)$ que
        elegiste.
\end{description}

% =============================================================================
\section{Cifrado RSA}
% =============================================================================

El cifrado RSA de esta actividad es \textbf{personalizado}: opera sobre
bloques numéricos formados a partir de grupos de tres caracteres.  Lee
esta sección con cuidado antes de empezar.

\subsection{Parámetros}

Elige los siguientes valores y verifica las condiciones indicadas:

\begin{enumerate}
  \item Dos primos distintos $p$ y $q$ tal que
        $10\,000 \le p,\, q \le 20\,000$.

  \item Un entero $e$ con $200 \le e \le 300$ tal que
        $\gcd\!\left(e,\; \phi(n)\right) = 1$,\; donde
        $\phi(n) = (p-1)(q-1)$.

  \item Calcula $n = p \cdot q$.

  \item Calcula $d$ tal que $e \cdot d \equiv 1 \pmod{\phi(n)}$.
        Este es tu exponente de decodificación.
\end{enumerate}

\begin{description}
  \item[Clave de codificación a reportar.] El par $(e,\, n)$.
  \item[Clave de decodificación a reportar.] El par $(d,\, n)$.
\end{description}

\subsection{Construcción de bloques}

Una vez que tienes tu frase en mayúsculas, sigue estos pasos:

\begin{enumerate}
  \item \textbf{Asigna códigos.}  Reemplaza cada carácter por su código
        de dos dígitos según la tabla.

  \item \textbf{Agrupa de derecha a izquierda} en grupos de tres
        caracteres.  Si la longitud del mensaje no es múltiplo de tres,
        el primer grupo (el de la izquierda) quedará incompleto; rellena
        con \textbf{espacios} (\texttt{40}) \textbf{por la izquierda}
        hasta completar tres caracteres.

  \item \textbf{Forma el valor numérico del bloque} concatenando los
        tres códigos de dos dígitos.  Si el bloque es
        $[c_1,\, c_2,\, c_3]$, entonces:
        \[
          M = c_1 \times 10\,000 + c_2 \times 100 + c_3
        \]
        El resultado es un número de hasta seis dígitos.

  \item \textbf{Cifra cada bloque:}
        \[
          C_i = M_i^{\,e} \bmod n
        \]
\end{enumerate}

\begin{ejemplo}[Construcción de bloques para la frase \texttt{HOLA}]

\textbf{Paso 1 — Asignar códigos:}

\begin{center}
\renewcommand{\arraystretch}{1.2}
\begin{tabular}{ccccc}
\toprule
H & O & L & A \\
\midrule
07 & 14 & 11 & 00 \\
\bottomrule
\end{tabular}
\end{center}

\textbf{Paso 2 — Agrupar de derecha a izquierda:}

El mensaje tiene 4 caracteres.  Como $4 = 1 \times 3 + 1$, el primer
bloque tendrá solo 1~carácter (\texttt{H}); se rellena con dos
espacios por la izquierda.

\textbf{Paso 3 — Valor numérico de cada bloque:}

\begin{center}
\renewcommand{\arraystretch}{1.3}
\begin{tabular}{clll}
\toprule
\textbf{Bloque} & \textbf{Caracteres} & \textbf{Códigos} & \textbf{Valor $M_i$} \\
\midrule
$M_1$ & [\texttt{esp}, \texttt{esp}, H] & [40, 40, 07] & 404\,007 \\
$M_2$ & [O, L, A]                       & [14, 11, 00] & 141\,100 \\
\bottomrule
\end{tabular}
\end{center}

\textbf{Paso 4 — Cifrar:}\quad
$C_1 = 404007^e \bmod n$,\quad $C_2 = 141100^e \bmod n$.

El texto cifrado se entrega como la secuencia de valores $C_i$
separados por comas: $C_1, C_2, \ldots$

\end{ejemplo}

\begin{advertencia}
  Verifica que $M_i < n$ para todo bloque.  Como $M_i \le 474\,747$
  y $n \ge 10\,000 \times 10\,001 = 100\,010\,000$, esta condición
  siempre se cumple con los parámetros de esta actividad.
\end{advertencia}

\subsection{Formato del texto cifrado RSA}

Entrega los valores $C_i$ como una \textbf{lista de enteros separados
por comas}, en el mismo orden en que aparecen los bloques en el
mensaje (de izquierda a derecha).  Ejemplo:
\[
  C_1,\; C_2,\; C_3,\; \ldots
\]

% =============================================================================
\section{Resumen de lo que debes subir}
% =============================================================================

\begin{center}
\renewcommand{\arraystretch}{1.4}
\begin{tabular}{llll}
\toprule
\textbf{Cifrado} & \textbf{Clave de codificación} &
\textbf{Clave de decodificación} & \textbf{Formato del cifrado} \\
\midrule
Desplazamiento & $k$ & —        & cadena de caracteres \\
Afín    & $(a,\, b)$ & —        & cadena de caracteres \\
RSA     & $(e,\, n)$ & $(d,\, n)$ & lista de enteros $C_i$ \\
\bottomrule
\end{tabular}
\end{center}

Para los tres cifrados también debes reportar el nombre de la
canción y artista, y la frase original en mayúsculas.\\[8pt]

\noindent\textbf{Acceso al Forms:}\quad
\url{https://forms.gle/9hBfr1R9P7mp9R1T9}

\end{document}
